% from ega4-8-fr.tex
\oldpage[IV-3]{28}
faisceau induit sur $S$ par tous les $\mathcal O_{S_\lambda}$~; en particulier
pour tout $s\in S$, l'anneau local $\mathcal O_s$ est \emph{le même} pour $S$
et pour tous les $S_\lambda$. Si $\mathcal N_\lambda$ est le Nilradical de
$\mathcal O_{S_\lambda}$, le Nilradical $\mathcal N$ de $\mathcal O_S$ a donc
en chaque point de $S$ la même fibre $\mathcal N_s$ (nilradical de
$\mathcal O_s$) que $u_\lambda^*(\mathcal N_\lambda)$ (induit sur $S$ par
$\mathcal N_\lambda$). L'hypothèse que $S$ est réduit entraîne par suite
$u_\lambda^*(\mathcal N_\lambda)=0$~; comme $\mathcal N_0$ est supposé de
type fini, il en est de même des $\mathcal N_\lambda$, et la conclusion résulte
donc de (8.5.7).
\end{proof}

\begin{corollary}[8.7.3]
\label{IV.8.7.3-fr}
Supposons vérifiée l'une des hypothèses suivantes~:
\begin{enumerate}
\item[\emph{a)}] $S_0$ est un préschéma noethérien et les morphismes
$u_{\lambda\mu}:S_\mu\to S_\lambda$ sont des immersions ouvertes.
\item[\emph{b)}] Les morphismes $u_{\lambda\mu}:S_\mu\to S_\lambda$ sont
fidèlement plats.
\end{enumerate}
Alors, pour que $S$ soit intègre, il faut et il suffit qu'il existe
$\lambda_0$ tel que $S_\lambda$ soit intègre pour $\lambda\geq\lambda_0$.
\end{corollary}

\begin{proof}
Dire qu'un préschéma est intègre signifie qu'il est à la fois réduit et
irréductible~; le corollaire résulte donc de (8.7.2) et de (8.4.3).
\end{proof}

\begin{remark}[8.7.4]
\label{IV.8.7.4-fr}
Si l'on ne fait pas d'hypothèse sur les $u_{\lambda\mu}$, il peut se faire que
$S$ soit intègre bien que tous les $S_\lambda$ soient non réduits et non
connexes, comme le montre l'exemple (8.4.4), où on prend l'anneau $A$ non
réduit.
\end{remark}

\subsection{Préschémas de présentation finie sur une limite projective de préschémas}
\label{subsection:IV.8.8-fr}

\begin{env}[8.8.1]
\label{IV.8.8.1-fr}
Conservant toujours les notations et hypothèses de (8.2.2), nous supposerons
donnés dans cette section deux $S_\alpha$-préschémas $X_\alpha$, $Y_\alpha$,
ce qui définit (8.2.5) deux systèmes projectifs de préschémas
$(X_\lambda,v_{\lambda\mu})$ et $(Y_\lambda,w_{\lambda\mu})$ en posant
$X_\lambda=X_\alpha\times_{S_\alpha}S_\lambda$,
$Y_\lambda=Y_\alpha\times_{S_\alpha}S_\lambda$,
$v_{\lambda\mu}=1_{X_\alpha}\times u_{\lambda\mu}$,
$w_{\lambda\mu}=1_{Y_\alpha}\times u_{\lambda\mu}$ (pour
$\alpha\leq\lambda\leq\mu$), dont les limites projectives sont respectivement
$X=X_\alpha\times_{S_\alpha}S$, $Y=Y_\alpha\times_{S_\alpha}S$, les
morphismes canoniques $v_\lambda:X\to X_\lambda$ et
$w_\lambda:Y\to Y_\lambda$ étant respectivement égaux à
$1_{X_\alpha}\times u_\lambda$ et $1_{Y_\alpha}\times u_\lambda$. Pour
$\alpha\leq\lambda\leq\mu$, on a une application canonique
$e_{\mu\lambda}:\operatorname{Hom}_{S_\lambda}(X_\lambda,Y_\lambda)\to
\operatorname{Hom}_{S_\mu}(X_\mu,Y_\mu)$, qui à tout $S_\lambda$-morphisme
$f_\lambda:X_\lambda\to Y_\lambda$ fait correspondre
$f_\mu=f_\lambda\times1_{S_\mu}:X_\lambda\times_{S_\lambda}S_\mu\to
Y_\lambda\times_{S_\lambda}S_\mu$, et il est clair que
$(\operatorname{Hom}_{S_\lambda}(X_\lambda,Y_\lambda),e_{\mu\lambda})$ est un
\emph{système inductif} d'ensembles. De même, on a une application canonique
$e_\lambda:\operatorname{Hom}_{S_\lambda}(X_\lambda,Y_\lambda)\to
\operatorname{Hom}_S(X,Y)$ qui à $f_\lambda$ fait correspondre
$f=f_\lambda\times1_S:X_\lambda\times_{S_\lambda}S\to
Y_\lambda\times_{S_\lambda}S$ et $(e_\lambda)$ est un système inductif
d'applications~; d'où, en passant à la limite inductive, une application
canonique, fonctorielle en $S_\alpha$, $X_\alpha$ et $Y_\alpha$~:
\begin{equation}
e:\varinjlim\operatorname{Hom}_{S_\lambda}(X_\lambda,Y_\lambda)
\longrightarrow\operatorname{Hom}_S(X,Y).
\tag{8.8.1.1}\label{IV.8.8.1.1-fr}
\end{equation}
\end{env}

\begin{theorem}[8.8.2]
\EGAReaderTarget{ega4.c04.thm.8-8-2}{IV.8.8.2-fr}
\begin{enumerate}
\item[(i)] Supposons $X_\alpha$ quasi-compact (resp. quasi-compact et
quasi-séparé), et $Y_\alpha$ localement de type fini (resp. localement de
présentation finie) sur $S_\alpha$. Alors l'application (8.8.1.1) est
injective (resp. bijective).
\item[(ii)] Supposons $S_0$ quasi-compact et quasi-séparé. Pour tout préschéma
$X$ de présentation finie sur $S$, il existe un $\lambda\in L$, un préschéma
$X_\lambda$ de présentation finie sur $S_\lambda$, et un $S$-isomorphisme
$X\xrightarrow{\sim}X_\lambda\times_{S_\lambda}S$.
\end{enumerate}
\end{theorem}

\begin{proof}
